\documentclass[11pt]{article}
\usepackage{amsmath, amssymb, amsthm}
\usepackage[margin=1in]{geometry}
\usepackage{hyperref}
\usepackage{graphicx}

\title{Chaos in the Lorenz System:\\ The Butterfly Effect and the Lyapunov Exponent}
\author{}
\date{}

\begin{document}
\maketitle

\section{The Lorenz System}

The Lorenz system is a set of three coupled, nonlinear ordinary differential
equations, originally derived by Edward Lorenz in 1963 as a drastically
simplified model of atmospheric convection:

\begin{align}
\dot{x} &= \sigma (y - x) \\
\dot{y} &= x(\rho - z) - y \\
\dot{z} &= xy - \beta z
\end{align}

where $x, y, z \in \mathbb{R}$ are the state variables and $\sigma, \rho,
\beta > 0$ are parameters. The classic chaotic regime uses
\[
\sigma = 10, \qquad \rho = 28, \qquad \beta = \frac{8}{3}.
\]

Despite being a smooth, deterministic, three-dimensional system, this choice
of parameters produces trajectories that never repeat, remain bounded within
a fractal set (the \emph{attractor}), and are extraordinarily sensitive to
initial conditions.

\section{Numerical Integration: 4th-Order Runge--Kutta}

Given a system $\dot{\mathbf{s}} = f(\mathbf{s})$ with state
$\mathbf{s} = (x, y, z)$, the classic RK4 method advances the state by a
time step $h$ via

\begin{align}
k_1 &= f(\mathbf{s}_n) \\
k_2 &= f\!\left(\mathbf{s}_n + \tfrac{h}{2} k_1\right) \\
k_3 &= f\!\left(\mathbf{s}_n + \tfrac{h}{2} k_2\right) \\
k_4 &= f(\mathbf{s}_n + h\, k_3) \\
\mathbf{s}_{n+1} &= \mathbf{s}_n + \frac{h}{6}\left(k_1 + 2k_2 + 2k_3 + k_4\right)
\end{align}

This is fourth-order accurate: the local truncation error per step is
$O(h^5)$, and the global error over a fixed time interval is $O(h^4)$. The
implementation uses $h = 0.005$, which is small enough to keep the
integration well within its stability region for the Lorenz parameters used
here, while remaining fast enough to simulate tens of thousands of steps in
a fraction of a second.

\section{Sensitive Dependence on Initial Conditions}

Let $\mathbf{s}_A(t)$ and $\mathbf{s}_B(t)$ be two trajectories of the same
system with initial separation
\[
\delta_0 = \lVert \mathbf{s}_B(0) - \mathbf{s}_A(0) \rVert = \varepsilon,
\]
where $\varepsilon$ is arbitrarily small (in this project, $\varepsilon =
10^{-8}$). For a chaotic system, the separation
$\delta(t) = \lVert \mathbf{s}_B(t) - \mathbf{s}_A(t) \rVert$ grows, on
average, \emph{exponentially}:
\[
\delta(t) \approx \delta_0\, e^{\lambda t}
\]
for as long as $\delta(t)$ remains small compared to the size of the
attractor. The constant $\lambda$ is the \textbf{largest Lyapunov
exponent}. Rearranging,
\[
\lambda \approx \frac{1}{t} \ln\!\left(\frac{\delta(t)}{\delta_0}\right).
\]

\begin{itemize}
  \item If $\lambda \le 0$: nearby trajectories converge or stay together — the system is \emph{not} chaotic.
  \item If $\lambda > 0$: nearby trajectories diverge exponentially — the system \emph{is} chaotic, and $1/\lambda$ is a natural estimate of the ``prediction horizon'': the timescale beyond which two initial conditions indistinguishable at measurement precision lead to unrelated outcomes.
\end{itemize}

For the classic Lorenz parameters, the widely cited reference value is
\[
\lambda \approx 0.905\ \text{(per unit time)}.
\]

\section{Benettin's Algorithm}

Naively integrating two trajectories $\varepsilon$ apart and measuring
$\delta(t)$ directly runs into a problem: $\delta(t)$ grows exponentially,
so within a short time it saturates at the size of the attractor itself
(here, tens of units) and the exponential-growth signal is lost in the
nonlinear folding of the attractor. Benettin's method (also called the
\emph{renormalization method}) fixes this:

\begin{enumerate}
  \item Integrate a reference trajectory $\mathbf{s}_{\text{ref}}(t)$ and a
        perturbed trajectory $\mathbf{s}_{\text{pert}}(t)$, starting a
        distance $\varepsilon$ apart, for a short interval of $N$ steps.
  \item Measure the new separation $d = \lVert \mathbf{s}_{\text{pert}} -
        \mathbf{s}_{\text{ref}} \rVert$.
  \item Accumulate $\ln(d / \varepsilon)$ into a running sum.
  \item \textbf{Rescale} $\mathbf{s}_{\text{pert}}$ back to distance exactly
        $\varepsilon$ from $\mathbf{s}_{\text{ref}}$, \emph{along the same
        direction of divergence}:
        \[
        \mathbf{s}_{\text{pert}} \leftarrow \mathbf{s}_{\text{ref}} +
        \frac{\varepsilon}{d}\left(\mathbf{s}_{\text{pert}} -
        \mathbf{s}_{\text{ref}}\right)
        \]
  \item Repeat for many renormalization cycles $r = 1, \dots, M$, each
        covering a time interval $N h$.
\end{enumerate}

After $M$ renormalizations covering total time $T = M \cdot N h$, the
largest Lyapunov exponent is estimated as
\[
\lambda \approx \frac{1}{T} \sum_{r=1}^{M} \ln\!\left(\frac{d_r}{\varepsilon}\right).
\]

This works because rescaling after every short interval keeps the
perturbation small enough to stay in the locally linear regime around the
reference trajectory (where the exponential-divergence approximation is
valid), while the accumulated log-growth still correctly captures the
average exponential rate over the whole attractor, including its folding
and stretching.

\subsection*{Parameters used in this project}
\[
\varepsilon = 10^{-8}, \quad N = 10 \text{ steps/renorm}, \quad
M = 4000 \text{ renorms}, \quad h = 0.005
\]
giving a total integration time of $T = 4000 \times 10 \times 0.005 = 200$
time units, over which the running estimate of $\lambda$ converges to
\[
\lambda \approx 0.9032,
\]
within $0.3\%$ of the reference value $0.905$.

\section{Summary}

\begin{center}
\begin{tabular}{ll}
\textbf{Quantity} & \textbf{Value / Method} \\
\hline
System & Lorenz equations, $\sigma=10,\ \rho=28,\ \beta=8/3$ \\
Integrator & RK4, step size $h=0.005$ \\
Chaos test & Two trajectories, $\varepsilon = 10^{-8}$ apart \\
Divergence & Exponential: $\delta(t) \approx \varepsilon\, e^{\lambda t}$ \\
Lyapunov exponent & $\lambda \approx 0.903$ (computed) vs.\ $0.905$ (reference) \\
Prediction horizon & $1/\lambda \approx 1.1$ time units \\
\end{tabular}
\end{center}

\end{document}
